% Generated draft: requires mathematical and source review.
\documentclass[11pt]{article}
\usepackage[margin=25mm]{geometry}
\usepackage{fontspec}
\setmainfont{Latin Modern Roman}
\usepackage{amsmath,amssymb,mathtools,unicode-math}
\setmathfont{Latin Modern Math}
\usepackage{enumitem,needspace,xurl,hyperref}
\setlength{\parindent}{0pt}
\setlength{\parskip}{5pt}
\setlength{\emergencystretch}{3em}
\newcommand{\N}{\mathbb{N}}
\newcommand{\Z}{\mathbb{Z}}
\newcommand{\Q}{\mathbb{Q}}
\newcommand{\R}{\mathbb{R}}
\newcommand{\C}{\mathbb{C}}
\newcommand{\F}{\mathbb{F}}
\newcommand{\A}{\mathbb{A}}
\newcommand{\PP}{\mathbb{P}}
\newcommand{\E}{\mathbb{E}}
\newcommand{\eps}{\varepsilon}
\newcommand{\dd}{\,\mathrm{d}}
\begin{document}
\section*{UnsolvedMath ID: 700326}
\textbf{Source problem code:} catalogue entry
\subsection{Definitions and mathematical statement}
A variety of groups is a class of groups defined by a set of group identities. Order varieties by inclusion. For a family \((\mathfrak V_\alpha)_{\alpha\in I}\), write \(\bigvee_{\alpha\in I}\mathfrak V_\alpha\) for its join: the smallest variety containing every \(\mathfrak V_\alpha\). The source's union notation is interpreted as this lattice join.

A representation
\[
\mathfrak A=\bigvee_{\alpha\in I}\mathfrak V_\alpha
\]
is called \emph{irreducible} in the source if
\[
\forall\alpha\in I,\qquad
\bigvee_{\beta\in I\setminus\{\alpha\}}\mathfrak V_\beta\ne\mathfrak A.
\]
A variety \(\mathfrak V\) cannot be decomposed into a join of two proper subvarieties if, for all varieties \(\mathfrak U,\mathfrak W\subseteq\mathfrak V\),
\[
\mathfrak V=\mathfrak U\vee\mathfrak W
\quad\Longrightarrow\quad
\mathfrak U=\mathfrak V\ \text{or}\ \mathfrak W=\mathfrak V.
\]

The question is whether every variety \(\mathfrak A\) admits an irreducible representation as above, with finitely or infinitely many constituent varieties \(\mathfrak V_\alpha\), each of which cannot be decomposed into a join of two proper subvarieties.

There is an unresolved boundary convention: the source does not explicitly specify whether the indexing set \(I\) may be empty. With the usual complete-lattice convention, the empty join is the variety \(\mathbf 1\) containing only trivial groups. Under this convention the empty representation of \(\mathbf 1\) satisfies the displayed irreducibility condition vacuously, whereas no nonempty representation of \(\mathbf 1\) satisfies it. Thus the admissibility of an empty family must be specified when interpreting the question for every variety.
\subsection{Short English statement}
Can every variety of groups be expressed as a finite or infinite lattice join in which no constituent can be omitted and each constituent cannot itself be expressed as a join of two proper subvarieties? The source leaves the empty-family convention unspecified, which affects the trivial variety.
\subsection{Sources}
\begin{itemize}
\item Local UnsolvedMath registry, permanent ID \(700326\), exact source problem code \texttt{catalogue entry}, title \emph{Kourovka Notebook Problem 5.44}. The supplied registry labels the problem open; this is a supplied status, not an independently verified mathematical or literature assessment.
\item \emph{The Kourovka Notebook (Unsolved Problems in Group Theory)}, Problem 5.44, originating in the fifth issue (1976), as cited by the supplied record through the twenty-first issue edition (2026), arXiv:1401.0300v48. The collection description names E. I. Khukhro and V. D. Mazurov as editors. This source citation is supplied and has not been independently checked.
\item The supplied proposer attribution begins A. Yu. Ol'shanski and ends in an incomplete TeX accent sequence. The registry gives the proposed year as 1976; the malformed name ending is not reconstructed here.
\end{itemize}
\end{document}
